\documentclass[11pt,a4paper]{article}
\usepackage[utf8]{inputenc}
\usepackage[T1]{fontenc}
\usepackage{lmodern}
\usepackage[margin=1in]{geometry}
\usepackage{amsmath,amssymb,amsfonts}
\usepackage{booktabs}
\usepackage{tabularx}
\usepackage{array}
\usepackage{tcolorbox}
\usepackage[hidelinks]{hyperref}

\title{\textbf{The Illuminated Room: Wave Propagation in Functional Infinity, Open-System Optical Condensation, and Cortical Temporal Binding in AEON-$T_{\mu\nu}$}}
\author{Lu\'is Pimenta Dam\'asio\\
\small Independent Researcher, Lisbon, Portugal\\
\small ORCID: \href{https://orcid.org/0009-0001-7606-7557}{0009-0001-7606-7557}}
\date{November 18, 2026}

\begin{document}

\maketitle

\begin{abstract}
An ordinary room illuminated by a single light bulb appears instantaneously lit to human perception. Yet from the standpoint of relativistic temporal metrology and open-system quantum foundations, this everyday scene compresses the entire operational architecture of AEON-$T_{\mu\nu}$ into a three-meter baseline. This paper traces the physical and ontological sequence of illumination across a $3\,\mathrm{m}$ room: from emission at the filament into uncoupled electromagnetic wave propagation within \emph{functional infinity} ($\Phi_\tau = 1$), through refractive phase-drag in transparent media ($v = c / n_{\mathrm{eff}}$), frequency-selective absorption and re-emission at colored surfaces, transverse phase filtering in polarizers, and geometric flux occlusion in shadows, to irreversible open-system condensation ($\Phi_\tau \to \Phi_{\mathrm{floor}}$) at retinal photoreceptors. Finally, this paper demonstrates why discrete retinal absorption events are experienced as a continuous, unified visual present only because thalamocortical phase-locking expands the cortical operational simultaneity horizon ($L_{\mathrm{sim}}^{\mathrm{cortex}} > L_{\mathrm{cortex}}$, $\mathcal{D}_{\mathrm{loc}} \to 1$), integrating discrete causal arrivals ($\Delta\tau_{\mathrm{sig}} \approx 20\,\mathrm{ns}$) across the biological integration window ($\sim 80\text{--}120\,\mathrm{ms}$).
\end{abstract}

\begin{tcolorbox}[colback=gray!6,colframe=gray!50,title=\textbf{Epistemic Status \& AI Collaboration Disclosure}]
\small
This paper presents a philosophical and theoretical synthesis grounded in established classical electrodynamics, physical optics, open quantum systems, and visual neuroscience. All quantitative optical constants, propagation delays, and neurobiological integration timescales rest on peer-reviewed literature, whereas the interpretive categories of \emph{functional infinity}, \emph{object-as-subject}, and \emph{diachronic identity} ($A_s^v$) belong to the speculative ontology of AEON-$T_{\mu\nu}$. AI assistance was utilized for mathematical typesetting, structural organization, and document preparation.
\end{tcolorbox}

\section{The Illusion of Instantaneous Illumination}

Switch on a lamp in the center of a $3\,\mathrm{m} \times 3\,\mathrm{m}$ room. To the observer standing beside the wall, the room appears to flood with light instantaneously and uniformly. In four-dimensional coordinate abstractions, this impression is reinforced by treating space as a simultaneous three-dimensional slice illuminated at $t = t_0$.

Within \textbf{AEON-$T_{\mu\nu}$} (developed across our \emph{Master Compendium (v2.0)} and \emph{AEON-$T_{\mu\nu}$ \& Philosophy of Time: Executive Manifesto and Front-Door Reader's Guide}), nothing in the room happens instantaneously. Every physical interaction unfolds at the finite causal pace of round-trip signaling lag ($\Delta\tau_{\mathrm{sig}}$) and open-system decoherence ($\Gamma_{\mathrm{tot}} = \Gamma_{\mathrm{th}} + \Gamma_{\mathrm{grav}}$):

\begin{itemize}
    \item \textbf{One-Way Optical Traversal ($L = 3\,\mathrm{m}$):} Electromagnetic radiation crosses the room in $\Delta t_{\mathrm{1-way}} = L / c \approx 10.0\,\mathrm{ns}$, corresponding to a round-trip closed-loop signaling lag of $\Delta\tau_{\mathrm{sig}} = 2L / c \approx 20.0\,\mathrm{ns}$.
    \item \textbf{Environmental Decoherence Rate ($\Gamma_{\mathrm{tot}}$):} In ambient room-temperature air ($T \approx 295\,\mathrm{K}$, $P \approx 1\,\mathrm{atm}$), collisional and thermal-bath scattering rates act as a continuous open-system environment that governs where waves remain phase-coherent versus where they condense into localized thermal events.
    \item \textbf{Human Cortical Integration Window ($\tau_{\mathrm{bind}} \approx 100\,\mathrm{ms}$):} Because the observer's thalamocortical simultaneity horizon ($L_{\mathrm{sim}}^{\mathrm{cortex}} > L_{\mathrm{cortex}}$) integrates sensory arrivals across roughly $10^7$ times the $10\,\mathrm{ns}$ optical transit time, the discrete causal buildup of illumination is smoothed into a single seamless perceptual ``Now.''
\end{itemize}

\section{Wave Propagation in Functional Infinity \& Refractive Phase-Drag}

Between thermal emission at the incandescent or LED source and final absorption at a wall or eye, the electromagnetic field does not travel as miniature solid pellets bouncing along classical trajectories. Rather, it propagates as a physically real, spatially extended wave manifold inhabiting \textbf{functional infinity} ($\Phi_\tau(\Omega) = 1$):

\begin{enumerate}
    \item \textbf{Vacuum \& Air Traversal ($\Phi_\tau = 1$):} Across the $3\,\mathrm{m}$ air gap ($n_{\mathrm{air}} \approx 1.00029$), coupling to sparse nitrogen and oxygen molecules is weak enough that phase coherence across optical wavelengths ($\lambda \approx 400\text{--}700\,\mathrm{nm}$) remains intact ($\Omega = \int \Gamma_{\mathrm{tot}}\,dt_{\mathrm{pre}} \ll 1$). The wave explores all geometrically accessible paths simultaneously without collapsing into localized particles mid-flight.
    \item \textbf{Refractive Phase-Drag in Glass and Water ($v = c / n_{\mathrm{eff}}$):} When the wavefront enters a window pane ($n_{\mathrm{glass}} \approx 1.50$) or a glass of water ($n_{\mathrm{water}} \approx 1.33$), coherent dipole polarization of the bound electronic lattice induces macroscopic phase-drag. The effective signal propagation speed drops to $v = c / n_{\mathrm{eff}}$, increasing the round-trip signaling lag to $\Delta\tau_{\mathrm{sig}} = 2\Delta\ell\,n_{\mathrm{eff}} / c$ while preserving wave coherence across transparent bulk media.
\end{enumerate}

\section{Reflection, Color, Polarization \& Shadow Occlusion}

As the wave manifold encounters macroscopic objects inside the room (sitting at conformal tier $\zeta = \log_{10}(L / \ell_P) \approx 34\text{--}35$), distinct open-system boundary interactions govern what survives in functional infinity versus what condenses thermally:

\begin{itemize}
    \item \textbf{Specular Reflection:} At a silvered mirror, conduction electrons oscillate coherently with the incident wavefront without dissipative thermalization, re-radiating a phase-locked wavefront that extends the causal path length ($\Delta\tau_{\mathrm{sig}}$) while keeping $\Phi_\tau \approx 1$ until detector arrival.
    \item \textbf{Color as Frequency-Selective Condensation:} A painted red wall acts as a frequency-selective open-system bath. Higher-frequency blue and green spectral components drive resonant electronic excitations that couple rapidly to lattice phonons ($\Omega \gg 1$), condensing irreversibly into local heat ($\Phi_\tau \to \Phi_{\mathrm{floor}}$), while lower-frequency red components scatter elastically back into the room.
    \item \textbf{Polarization Filtering:} Passing the beam through a dichroic polarizer projects transverse field oscillations onto a single axis, absorbing the orthogonal component into molecular vibrational modes ($\Phi_\tau \to \Phi_{\mathrm{floor}}$) while transmitting the aligned component coherently.
    \item \textbf{Shadow Occlusion:} An opaque chair placed $1.5\,\mathrm{m}$ from the lamp intercepts the wavefront, condensing incident energy into surface heat and casting a geometric shadow cone across the far wall. Within the umbra, direct radiative flux drops to zero, leaving only multiply scattered ambient waves with longer causal traversal histories.
\end{itemize}

\section{Retinal Condensation, Cortical Simultaneity \& Methodological Note}

Vision occurs only when a portion of the room's reflected wave manifold enters the pupil and strikes a retinal rhodopsin or photopsin chromophore. At that molecular locus, open-system coupling drives $\Omega \gg 1$, condensing the extended wave manifold irreversibly into a localized photo-isomerization event ($\Phi_\tau \to \Phi_{\mathrm{floor}}$).

\subsection{From Discrete Photochemical Events to Unified Visual Presence}
Individual retinal isomerizations occur at staggered proper times across the photoreceptor mosaic, separated by nanoseconds of optical path difference and milliseconds of synaptic conduction delay. They are unified into a single coherent visual scene only because the observer is a metabolically phase-locked \emph{object-as-subject} whose thalamocortical oscillations sustain $L_{\mathrm{sim}}^{\mathrm{cortex}} > L_{\mathrm{cortex}}$ ($\mathcal{D}_{\mathrm{loc}} \to 1$). Across the $\sim 100\,\mathrm{ms}$ cortical integration window (and its brief retinal/cortical afterimage persistence), discrete physical arrivals are bound into a continuous diachronic identity ($A_s^v$).

\begin{tcolorbox}[colback=blue!4,colframe=blue!40,title=\textbf{Methodological Footnote: Apparatus-Dependent Open-System Condensation}]
\small
A crucial operational lesson follows from treating detection as open-system wave condensation ($\Phi_\tau = 1 \to \Phi_{\mathrm{floor}}$) rather than passive readout of pre-existing point properties. What an instrument registers depends jointly on the incoming wave manifold and the specific dissipative coupling ($\Gamma_{\mathrm{tot}}$), geometry, and integration horizon ($\mathcal{L}_{\mathrm{sim}}$) of the detecting apparatus. While benign in an illuminated room, this principle becomes acute in precision experimental physics: for example, persistent $5\sigma+$ discrepancies between independent electron--positron collider detectors (such as KLOE, BaBar, and CMD-3) measuring the ``same'' hadronic cross-section ($e^+e^- \to \pi^+\pi^-$) underscore how distinct detector geometries, magnetic fields, and radiative-return timing windows sample open-system condensation differently from idealized closed-box lattice QCD calculations.
\end{tcolorbox}

\section{Comparative Positioning Across Four Traditions}

\begin{table}[h!]
\centering
\scriptsize
\begin{tabularx}{\textwidth}{>{\bfseries}p{2.5cm} X X X}
\toprule
Framework & Light in Transit & Surface \& Retinal Interaction & Unified Visual ``Now'' \\
\midrule
AEON-$T_{\mu\nu}$ (This Paper) & Real wave manifold in functional infinity ($\Phi_\tau = 1$) with finite lag $\Delta\tau_{\mathrm{sig}}$ & Open-system thermal/photochemical condensation ($\Phi_\tau \to \Phi_{\mathrm{floor}}$) & Cortical operational simultaneity ($L_{\mathrm{sim}}^{\mathrm{cortex}} > L_{\mathrm{cortex}}$) \\
\addlinespace
Classical Maxwellian Electrodynamics & Continuous classical field propagating at $c / n_{\mathrm{eff}}$ & Continuous energy absorption via Poynting flux $\mathbf{S} = \mathbf{E} \times \mathbf{H}$ & Left to external psychophysics; no observer boundary \\
\addlinespace
Textbook Copenhagen Quantum Mechanics & Abstract probability amplitude $\psi(\mathbf{r},t)$ & Instantaneous non-unitary measurement collapse & Ambiguous observer cut (von Neumann chain) \\
\addlinespace
4D Block Universe Eternalism & Static null worldlines connecting emission and absorption events & Fixed 4D intersection vertices & Psychological illusion along worldline slices \\
\bottomrule
\end{tabularx}
\caption{Comparison of optical propagation, detection, and perceptual binding across four frameworks.}
\end{table}

\section{Four Explicit Structural Limitations}

\begin{enumerate}
    \item \textbf{Quantum-Gravity Cutoff ($\Xi_{\mathrm{QG}} < 1$):} Optical propagation in a room sits deep within the weak-field semiclassical regime ($\Xi_{\mathrm{QG}} \ll 10^{-30}$); this paper makes no claims regarding Planck-scale photon dispersion or quantum-gravitational foam.
    \item \textbf{Macroscopic Effective Refractive Index:} Treating glass or water via a scalar index $n_{\mathrm{eff}}$ is an effective macroscopic approximation valid far from sharp atomic absorption resonances and strong nonlinearities.
    \item \textbf{Hard Problem of Color Qualia:} While frequency-selective condensation and cortical phase-locking explain the structural timing and unity of color perception, this paper does not deduce the intrinsic subjective feel of phenomenal redness from optical equations.
    \item \textbf{Speculative Ontological Vocabulary:} Established electrodynamics and visual neuroscience account for all quantitative numbers herein; the categories of \emph{functional infinity}, \emph{object-as-subject}, and \emph{diachronic identity} ($A_s^v$) represent a philosophical synthesis open to empirical refinement.
\end{enumerate}

\section{Conclusion}

The illuminated room demonstrates that the core principles of \textbf{AEON-$T_{\mu\nu}$} are not confined to cosmological horizons or precision atomic clock laboratories. Even in a three-meter room, physical reality unfolds at a finite causal pace ($\Delta\tau_{\mathrm{sig}} \approx 20\,\mathrm{ns}$): light traverses space as an uncollapsed wave in functional infinity ($\Phi_\tau = 1$), condenses irreversibly through open-system interactions at surfaces and retinas ($\Phi_\tau \to \Phi_{\mathrm{floor}}$), and becomes a unified visual present only inside a phase-locked conscious subject ($L_{\mathrm{sim}}^{\mathrm{cortex}} > L_{\mathrm{cortex}}$).

\begin{thebibliography}{18}
\bibitem{DamasioCompendium1} Dam\'asio, L.~P. (2026). \emph{Master Compendium --- Foundations of Temporal Metrology and Operational Quantum Mechanics (v1.0)}. Zenodo. \url{https://doi.org/10.5281/zenodo.24163322}
\bibitem{DamasioCompendium2} Dam\'asio, L.~P. (2026). \emph{Master Compendium --- Foundations of Temporal Metrology and Operational Quantum Mechanics (v2.0)}. Zenodo. \url{https://doi.org/10.5281/zenodo.24190229}
\bibitem{DamasioObservatory} Dam\'asio, L.~P. (2026). \emph{AEON-T14.html: Interactive 6-Panel Computational Observatory across 120 Orders of Magnitude}. Zenodo. \url{https://doi.org/10.5281/zenodo.24244501}
\bibitem{DamasioManifesto} Dam\'asio, L.~P. (2026). \emph{AEON-$T_{\mu\nu}$ \& Philosophy of Time: Executive Manifesto and Front-Door Reader's Guide}. Zenodo. \url{https://doi.org/10.5281/zenodo.24339158}
\bibitem{DamasioNeuro} Dam\'asio, L.~P. (2026). \emph{Coherence-Bounded Temporal Binding: Cortical Operational Simultaneity, Object-as-Subject Interiority, and the Epigenetics of Consciousness in AEON-$T_{\mu\nu}$}. Zenodo. \url{https://doi.org/10.5281/zenodo.24339942}
\bibitem{DamasioBio} Dam\'asio, L.~P. (2026). \emph{Metabolic Coherence and Diachronic Identity: Biological Operational Simultaneity, Nested Organismic Clocks, and Thermodynamic Persistence in AEON-$T_{\mu\nu}$}. Zenodo. \url{https://doi.org/10.5281/zenodo.24345945}
\bibitem{DamasioGeodesy} Dam\'asio, L.~P. (2026). \emph{Coherence-Bounded Optical Geodesy: Relativistic Redshift, Operational Simultaneity Horizons, and Null Interferometric Visibility in Terrestrial and Space-Based Clock Networks}. Zenodo. \url{https://doi.org/10.5281/zenodo.24346852}
\bibitem{DamasioSlow} Dam\'asio, L.~P. (2026). \emph{Slow Metaphysics: Objective Temporal Cadence, Diachronic Identity ($A_s^v$), and Functional Infinity in AEON-$T_{\mu\nu}$}. Zenodo. \url{https://doi.org/10.5281/zenodo.24351508}
\bibitem{Maxwell1865} Maxwell, J.~C. (1865). A dynamical theory of the electromagnetic field. \emph{Philosophical Transactions of the Royal Society of London}, 155, 459--512.
\bibitem{BornWolf} Born, M., \& Wolf, E. (1999). \emph{Principles of Optics} (7th ed.). Cambridge University Press.
\bibitem{Glauber1963} Glauber, R.~J. (1963). The quantum theory of optical coherence. \emph{Physical Review}, 130(6), 2529--2539.
\bibitem{Zurek2003} Zurek, W.~H. (2003). Decoherence, einselection, and the quantum origins of the classical. \emph{Reviews of Modern Physics}, 75(3), 715--775.
\bibitem{JoosZeh} Joos, E., et al. (2003). \emph{Decoherence and the Appearance of a Classical World in Quantum Theory} (2nd ed.). Springer.
\bibitem{FeynmanQED} Feynman, R.~P. (1985). \emph{QED: The Strange Theory of Light and Matter}. Princeton University Press.
\bibitem{Wald1978} Wald, G. (1968). The molecular basis of visual excitation. \emph{Nature}, 219(5156), 800--807.
\bibitem{Fries2015} Fries, P. (2015). Rhythms for cognition: Communication through coherence. \emph{Neuron}, 88(1), 220--235.
\bibitem{Massimini2013} Casali, A.~G., et al. (2013). A theoretically based index of consciousness independent of sensory processing and behavior. \emph{Science Translational Medicine}, 5(198), 198ra105.
\bibitem{CMD3_2024} Ignatov, F.~V., et al. (CMD-3 Collaboration). (2024). Measurement of the $e^+e^- \to \pi^+\pi^-$ cross section from threshold to $1.2\,\mathrm{GeV}$ with the CMD-3 detector. \emph{Physical Review D}, 109(11), 112002.
\end{thebibliography}

\end{document}
